\documentclass[12 pt]{extarticle}

	\usepackage[frenchb]{babel}
	\usepackage[utf8]{inputenc}  
	\usepackage[T1]{fontenc}
	\usepackage{amssymb}
	\usepackage[mathscr]{euscript}
	\usepackage{stmaryrd}
	\usepackage{amsmath}
	\usepackage{tikz}
	\usepackage[all,cmtip]{xy}
	\usepackage{amsthm}
	\usepackage{varioref}
	\usepackage{geometry}
	\geometry{a4paper}
	\usepackage{lmodern}
	\usepackage{multicol}
	\usepackage{hyperref}
	\usepackage{array}
	 \usepackage{fancyhdr}
%\renewcommand{\theenumi}{\alph{enumi})}
	\pagestyle{fancy}
	\theoremstyle{plain}
	\fancyfoot[C]{} 
	\fancyhead[L]{Contrôle}
	\fancyhead[R]{29 septembre 2025}\geometry{
 a4paper,
 total={170mm,257mm},
 left=10mm,right=10mm,
 top=5mm, bottom = 10 mm
 }
	
	
	\title{Interrogation Chapitre 2}
	\date{}
	\begin{document}

%\begin{center}{\Large Interrogation}\\ 
% \end{center}
\subsection*{Exercice 1 (6 points)}
\begin{enumerate}
\item Calculer $\frac{12}{10} + \frac{94}{100}$.
\item Calculer $3,13 + \frac{54}{10}$.
\item Compléter avec le symbole correct parmi $>$, $<$ et $=$ : 
\[ 1,12 \ldots 1,14 \phantom{miaoumiaou}
1,24 \ldots 1,3 \phantom{miaoumiaou}
\]
\item Écrire en notation mathématique la phrase suivante : « Le point C appartient à la demi-droite d'origine O passant par B »
\end{enumerate}
\subsection*{Exercice 2 (6 points)}
\includegraphics[scale=.5]{Figure}\\
On se donne la figure $ABCDEF$ ci-dessus. 
\begin{enumerate}
\item Que dire des angles $C$, $D$, $E$ et $F$ ? 
\item Démontrer que les droites $(DE)$ et $(CF)$ sont parallèles. 
\item En déduire que si l'on prolongeait la droite $(AD)$, elle rencontrerait $(CF)$ avec un angle droit.
\end{enumerate}

\subsection*{Exercice 3 (8 points)}
On considère la figure suivante où les droites $(BC)$ et $(AD)$ sont parallèles, et où les angles en $B$ et $D$ sont droits. \\
\begin{tikzpicture}
\draw (0,0) rectangle (5,3);
\draw (-.2, -.2) node{A}
(-.2, 3.2) node{B} 
(5.2, 3.2) node{C}
(5.2, -.2) node{D};
\end{tikzpicture}
\begin{enumerate}
\item Recopier et compléter avec le symbole adapté : $(AB) \ldots (BC)$
\item Recopier et compléter avec le symbole adapté : $(AD) \ldots (DC)$
\item Recopier et compléter avec le symbole adapté : $(AB) \ldots (DC)$
\item Démontrer que les droites $(AB)$ et $(AD)$ sont perpendiculaires.
\item Démontrer que les droites $(BC)$ et $(CD)$ sont perpendiculaires. 
\item Que dire du quadrilatère $ABCD$ ?

\end{enumerate}

\newpage

\subsection*{Exercice 1 (6 points)}
\begin{enumerate}
\item Calculer $\frac{59}{10} + \frac{12}{100}$.
\item Calculer $8,13 + \frac{54}{10}$.
\item Compléter avec le symbole correct parmi $>$, $<$ et $=$ : 
\[ 1,22 \ldots 1,14 \phantom{miaoumiaou}
1,2 \ldots 1,19 \phantom{miaoumiaou}
\]
\item Écrire en notation mathématique la phrase suivante : « Le point A appartient à la demi-droite d'origine M passant par N »
\end{enumerate}
\subsection*{Exercice 2 (6 points)}
\includegraphics[scale=.5]{Figure}\\
On se donne la figure $ABCDEF$ ci-dessus. 
\begin{enumerate}
\item Que dire des angles $C$, $D$, $E$ et $F$ ? 
\item Démontrer que les droites $(DE)$ et $(CF)$ sont parallèles. 
\item En déduire que si l'on prolongeait la droite $(AD)$, elle rencontrerait $(CF)$ avec un angle droit.
\end{enumerate}

\subsection*{Exercice 3 (8 points)}
On considère la figure suivante où les droites $(BC)$ et $(AD)$ sont parallèles, et où les angles en $B$ et $D$ sont droits. \\
\begin{tikzpicture}
\draw (0,0) rectangle (5,3);
\draw (-.2, -.2) node{A}
(-.2, 3.2) node{B} 
(5.2, 3.2) node{C}
(5.2, -.2) node{D};
\end{tikzpicture}
\begin{enumerate}
\item Recopier et compléter avec le symbole adapté : $(AB) \ldots (DC)$
\item Recopier et compléter avec le symbole adapté : $(AB) \ldots (BC)$
\item Recopier et compléter avec le symbole adapté : $(AD) \ldots (DC)$
\item Démontrer que les droites $(AB)$ et $(AD)$ sont perpendiculaires.
\item Démontrer que les droites $(BC)$ et $(CD)$ sont perpendiculaires. 
\item Que dire du quadrilatère $ABCD$ ?

\end{enumerate}
 
 	\end{document}
